% ============================================================
%  FireRisk Mediterranean — Technical Report
%  Deep Learning for Wildfire Risk Prediction
% ============================================================
\documentclass[12pt,a4paper]{report}

% ── Encoding & fonts ──────────────────────────────────────
\usepackage[utf8]{inputenc}
\usepackage[T1]{fontenc}
\usepackage[french]{babel}
\usepackage{lmodern}
\usepackage{microtype}

% ── Geometry & spacing ────────────────────────────────────
\usepackage[top=2.5cm,bottom=2.5cm,left=2.8cm,right=2.8cm]{geometry}
\usepackage{setspace}
\onehalfspacing
\setlength{\parskip}{6pt}
\setlength{\parindent}{0pt}

% ── Colors (Streamlit palette) ────────────────────────────
\usepackage[table,dvipsnames]{xcolor}
\definecolor{fireOrange}{HTML}{FF4500}
\definecolor{fireLight}{HTML}{FF8C42}
\definecolor{firePale}{HTML}{FF6B35}
\definecolor{darkBg}{HTML}{0F1420}
\definecolor{cardBg}{HTML}{1A1F2E}
\definecolor{lstmBlue}{HTML}{42A5F5}
\definecolor{tcnRed}{HTML}{EF5350}
\definecolor{riskGreen}{HTML}{4CAF50}
\definecolor{riskYellow}{HTML}{FFC107}
\definecolor{riskOrange}{HTML}{FF9800}
\definecolor{riskRed}{HTML}{F44336}
\definecolor{accentGray}{HTML}{AAAAAA}
\definecolor{lightGray}{HTML}{E0E0E0}
\definecolor{deepGreen}{HTML}{1B5E20}

% ── Graphics & figures ────────────────────────────────────
\usepackage{graphicx}
\usepackage{float}
\usepackage{wrapfig}
\usepackage{subcaption}
\usepackage{tikz}
\usetikzlibrary{shapes.geometric,arrows.meta,positioning,fit,backgrounds,shadows}

% ── Tables ────────────────────────────────────────────────
\usepackage{booktabs}
\usepackage{tabularx}
\usepackage{multirow}
\usepackage{array}
\usepackage{colortbl}

% ── Mathematics ───────────────────────────────────────────
\usepackage{amsmath,amssymb,amsthm}

% ── Code listings ─────────────────────────────────────────
\usepackage{listings}
\lstset{
  backgroundcolor=\color{cardBg},
  basicstyle=\ttfamily\small\color{lightGray},
  keywordstyle=\color{fireLight}\bfseries,
  commentstyle=\color{accentGray}\itshape,
  stringstyle=\color{riskGreen},
  numberstyle=\tiny\color{accentGray},
  frame=single,
  rulecolor=\color{fireOrange!40},
  numbers=left,
  breaklines=true,
  tabsize=4,
  language=Python,
}

% ── Hyperlinks ────────────────────────────────────────────
\usepackage[colorlinks=true,
            linkcolor=fireOrange,
            citecolor=lstmBlue,
            urlcolor=firePale,
            bookmarksopen=true]{hyperref}

% ── Headers / footers ─────────────────────────────────────
\usepackage{fancyhdr}
\pagestyle{fancy}
\fancyhf{}
\fancyhead[L]{\small\color{fireOrange}\textbf{FireRisk — Méditerranée}}
\fancyhead[R]{\small\color{accentGray}Prédiction du risque d'incendie par Deep Learning}
\fancyfoot[C]{\small\color{accentGray}\thepage}
\renewcommand{\headrulewidth}{0.4pt}
\renewcommand{\headrule}{\hbox to\headwidth{\color{fireOrange}\leaders\hrule height \headrulewidth\hfill}}

% ── Chapter / section styles ──────────────────────────────
\usepackage{titlesec}
\titleformat{\chapter}[block]
  {\normalfont\Huge\bfseries\color{fireOrange}}
  {\thechapter.}{1em}{}
  [\vspace{4pt}{\color{fireOrange!40}\titlerule[2pt]}]

\titleformat{\section}
  {\normalfont\Large\bfseries\color{firePale}}
  {\thesection}{1em}{}

\titleformat{\subsection}
  {\normalfont\large\bfseries\color{fireLight}}
  {\thesubsection}{1em}{}

\titleformat{\subsubsection}
  {\normalfont\normalsize\bfseries\color{accentGray}}
  {\thesubsubsection}{1em}{}

% ── Custom boxes ──────────────────────────────────────────
\usepackage[skins,breakable]{tcolorbox}

\newtcolorbox{firebox}[1][]{
  enhanced,colback=fireOrange!8,colframe=fireOrange,
  arc=6pt,boxrule=1pt,fonttitle=\bfseries\color{fireOrange},
  #1
}
\newtcolorbox{lstmbox}[1][]{
  enhanced,colback=lstmBlue!8,colframe=lstmBlue,
  arc=6pt,boxrule=1pt,fonttitle=\bfseries\color{lstmBlue},
  #1
}
\newtcolorbox{tcnbox}[1][]{
  enhanced,colback=tcnRed!8,colframe=tcnRed,
  arc=6pt,boxrule=1pt,fonttitle=\bfseries\color{tcnRed},
  #1
}
\newtcolorbox{infobox}[1][]{
  enhanced,colback=cardBg,colframe=accentGray!50,
  arc=6pt,boxrule=0.5pt,#1
}

% ── Misc ──────────────────────────────────────────────────
\usepackage{enumitem}
\usepackage{soul}
\usepackage{siunitx}
\usepackage{algorithm}
\usepackage{algpseudocode}
\usepackage{caption}
\captionsetup{labelfont={bf,color=fireOrange},font={small}}

% ── Bibliography ──────────────────────────────────────────
\usepackage[backend=biber,style=numeric,sorting=none]{biblatex}

% ============================================================
%  DOCUMENT
% ============================================================
\begin{document}

% ──────────────────────────────────────────────────────────
%  TITLE PAGE
% ──────────────────────────────────────────────────────────
\begin{titlepage}
  \begin{tikzpicture}[remember picture,overlay]
    % Background gradient rectangle
    \fill[darkBg] (current page.south west) rectangle (current page.north east);
    % Top fire accent bar
    \fill[fireOrange] (current page.north west)
      rectangle ([yshift=-8pt]current page.north east);
    % Orange glow circle top-right
    \fill[fireOrange,opacity=0.06] ([xshift=-3cm,yshift=-4cm]current page.north east)
      circle (7cm);
    % Bottom accent bar
    \fill[fireOrange!40,opacity=0.3] (current page.south west)
      rectangle ([yshift=4pt]current page.south east);
  \end{tikzpicture}

  \vspace*{3cm}
  \begin{center}
    % Icon
    {\fontsize{64}{72}\selectfont 🔥}\\[1cm]

    {\fontsize{30}{36}\selectfont\bfseries\color{fireOrange}
      FireRisk — Bassin Méditerranéen}\\[0.6cm]

    {\fontsize{16}{20}\selectfont\color{firePale}
      Prédiction du risque d'incendie de forêt par Deep Learning}\\[0.4cm]

    {\large\color{lightGray}
      LSTM $\cdot$ TCN $\cdot$ ERA5 $\cdot$ MODIS NDVI $\cdot$ NASA FIRMS}\\[2cm]

    \begin{tikzpicture}
      \draw[fireOrange,line width=1.5pt] (0,0) -- (10,0);
    \end{tikzpicture}\\[1.2cm]

    % ── Établissement & filière ──────────────────────────
    {\fontsize{13}{16}\selectfont\bfseries\color{fireOrange}
      Institut National des Postes et Télécommunications — INPT}\\[0.3cm]
    {\large\color{firePale}
      Filière : \textbf{DATA} \quad$|$\quad Module : \textbf{Deep Learning}}\\[1.2cm]

    \begin{tikzpicture}
      \draw[fireOrange!40,line width=0.5pt] (0,0) -- (10,0);
    \end{tikzpicture}\\[1cm]

    % ── Auteurs ─────────────────────────────────────────
    {\large\color{lightGray}\textbf{Réalisé par :}}\\[0.4cm]
    \begin{tikzpicture}
      \node[draw=fireOrange,fill=fireOrange!10,rounded corners=8pt,
            inner xsep=20pt,inner ysep=10pt,line width=1pt] {%
        \begin{tabular}{c}
          {\large\bfseries\color{fireOrange} Ilyass Aamoum}\\[4pt]
          {\large\bfseries\color{fireOrange} Chihab Chouai}
        \end{tabular}
      };
    \end{tikzpicture}\\[1cm]

    % ── Encadrant ───────────────────────────────────────
    {\large\color{lightGray}\textbf{Encadrant :}}\\[0.3cm]
    {\large\bfseries\color{firePale} Mr.\ Tarik Fissaa}\\[1cm]

    \begin{tikzpicture}
      \draw[fireOrange!30,line width=0.4pt] (0,0) -- (10,0);
    \end{tikzpicture}\\[0.6cm]

    {\color{accentGray}
      Période d'étude : 2022–2024\\
      Entraînement : 2022 \quad|\quad Test : 2023 (données jamais vues)\\[0.8cm]
      \today}
  \end{center}
\end{titlepage}

% ──────────────────────────────────────────────────────────
%  TABLE OF CONTENTS
% ──────────────────────────────────────────────────────────
\tableofcontents
\clearpage

% ──────────────────────────────────────────────────────────
%  ABSTRACT
% ──────────────────────────────────────────────────────────
\chapter*{Résumé}
\addcontentsline{toc}{chapter}{Résumé}

\begin{firebox}[title=Résumé du projet]
Le bassin méditerranéen est l'une des régions du globe les plus vulnérables aux incendies
de forêt. La conjonction d'étés chauds et secs, de végétation arbustive dense, et de vents
locaux forts (Mistral, Tramontane, Sirocco) crée régulièrement des conditions
météorologiques extrêmes favorisant la propagation rapide des feux.
Ce rapport présente \textbf{FireRisk}, un système de prédiction du risque d'incendie
fondé sur deux architectures de deep learning — \textcolor{lstmBlue}{\textbf{LSTM}}
(Long Short-Term Memory) et \textcolor{tcnRed}{\textbf{TCN}} (Temporal Convolutional
Network) — entraînées sur des données multi-sources couvrant l'intégralité du bassin
méditerranéen à une résolution de $0.25°$ ($\approx 28$ km).
\end{firebox}

\vspace{0.5cm}

Le système intègre trois sources de données principales :
\begin{itemize}[leftmargin=*,itemsep=2pt]
  \item \textbf{ERA5} (ECMWF) : réanalyse atmosphérique journalière (température,
        humidité, vent, précipitations) ;
  \item \textbf{MODIS NDVI/EVI} (NASA Terra/Aqua) : indice de végétation à résolution
        temporelle de 16 jours, projeté sur la grille d'étude ;
  \item \textbf{NASA FIRMS} : détections thermiques actives (hotspots) et puissance
        radiative du feu (FRP) en temps quasi-réel.
\end{itemize}

Les modèles consomment des \textbf{fenêtres glissantes de 15 jours} (8 features par
cellule) et prédisent un score de risque continu dans $[0,1]$ pour les 7 jours suivants.
Le protocole d'évaluation repose sur un \textbf{split temporel strict} : entraînement sur
2022, validation sur décembre 2022, test sur l'intégralité de l'année 2023 — données
jamais vues pendant l'entraînement.

Une \textbf{interface web interactive} (Streamlit) permet de visualiser les prédictions
ponctuelles, les cartes de risque sur toute la Méditerranée, et les courbes d'évaluation
des deux modèles.

\clearpage

% ──────────────────────────────────────────────────────────
%  CHAPTER 1 : INTRODUCTION
% ──────────────────────────────────────────────────────────
\chapter{Introduction}

\section{Contexte et motivation}

Les incendies de forêt constituent l'une des catastrophes naturelles les plus dévastatrices
du pourtour méditerranéen. Chaque été, des millions d'hectares de végétation sont consumés,
provoquant des pertes humaines, des destructions d'habitats, et des émissions massives de
CO\textsubscript{2}. L'été 2022 a été particulièrement meurtrier en France (incendies de
Gironde, $>20\,000$ ha brûlés), en Espagne, au Portugal et en Grèce. En 2023, les incendies
d'Evros (nord-est de la Grèce) ont été qualifiés de pires feux de forêt jamais enregistrés
en Europe, avec $>80\,000$ ha dévastés.

Face à l'intensification de ces événements sous l'effet du changement climatique, la
\textbf{prédiction précoce du risque} devient une priorité opérationnelle pour les agences
de protection civile, les services forestiers, et les assureurs. Les systèmes d'alerte
classiques (indices météorologiques comme le Feux de Forêt Weather Index) présentent des
limitations importantes : ils ignorent l'état dynamique de la végétation, l'historique
récent des détections et la mémoire à long terme des conditions atmosphériques.

\section{Problématique}

La prédiction du risque d'incendie est un problème \textbf{spatio-temporel complexe} car :
\begin{enumerate}[leftmargin=*]
  \item Le risque est fondamentalement \emph{non-stationnaire} : il évolue avec la saison,
        le climat, et la dynamique de la végétation ;
  \item Il existe une forte \emph{dépendance temporelle} : les jours précédant un feu sont
        caractérisés par une accumulation de stress thermique et hydrique ;
  \item La distribution des étiquettes est \emph{fortement déséquilibrée} : les cellules
        à risque élevé sont minoritaires, rendant l'apprentissage difficile.
\end{enumerate}

\section{Contributions du projet}

Ce projet apporte les contributions suivantes :

\begin{firebox}[title=Contributions principales]
\begin{enumerate}[leftmargin=*,itemsep=4pt]
  \item \textbf{Pipeline de données} complet intégrant ERA5, MODIS NDVI et NASA FIRMS
        sur une grille $0.25°$ couvrant le bassin méditerranéen (2022–2024) ;
  \item \textbf{Deux architectures deep learning} (LSTM bi-couche et TCN à dilations
        exponentielles) entraînées avec une perte pondérée adaptée au déséquilibre
        de classes ;
  \item \textbf{Protocole d'évaluation rigoureux} avec split temporel strict assurant
        l'absence de fuite d'information entre train et test ;
  \item \textbf{Interface web interactive} permettant la prédiction en temps réel et
        la visualisation de cartes de risque sur l'ensemble de la zone d'étude ;
  \item \textbf{Moteur d'inférence} optimisé supportant la prédiction par lot sur la
        grille complète ($49 \times 133 = 6\,517$ cellules).
\end{enumerate}
\end{firebox}

\section{Organisation du rapport}

Le rapport est organisé comme suit : le chapitre~\ref{chap:data} décrit les sources de
données et leur traitement. Le chapitre~\ref{chap:features} présente l'ingénierie des
features. Le chapitre~\ref{chap:models} décrit les architectures des modèles. Le
chapitre~\ref{chap:training} détaille le protocole d'entraînement. Le
chapitre~\ref{chap:eval} présente les résultats d'évaluation. Le
chapitre~\ref{chap:app} décrit l'interface web. Le
chapitre~\ref{chap:conclusion} conclut et propose des perspectives.

\clearpage

% ──────────────────────────────────────────────────────────
%  CHAPTER 2 : ZONE D'ÉTUDE ET DONNÉES
% ──────────────────────────────────────────────────────────
\chapter{Zone d'étude et Sources de données}
\label{chap:data}

\section{Zone d'étude : le bassin méditerranéen}

La zone d'étude couvre le bassin méditerranéen dans son intégralité :
\begin{equation}
  \mathcal{D} = \{(\phi, \lambda) \mid \phi \in [35°\text{N},\, 47°\text{N}],\;
  \lambda \in [5°\text{W},\, 28°\text{E}]\}
\end{equation}

\begin{center}
\begin{tabular}{lcc}
\toprule
\rowcolor{fireOrange!15}
\textbf{Paramètre} & \textbf{Valeur} & \textbf{Description} \\
\midrule
Latitude min/max    & $35.0°$N -- $47.0°$N  & 12° de largeur \\
Longitude min/max   & $5.0°$W -- $28.0°$E   & 33° de longueur \\
Résolution spatiale & $0.25°$ ($\approx 28$ km) & Grille régulière \\
Nombre de cellules lat & $N_\text{lat} = 49$    & $\lfloor 12/0.25 \rfloor + 1$ \\
Nombre de cellules lon & $N_\text{lon} = 133$   & $\lfloor 33/0.25 \rfloor + 1$ \\
Cellules totales    & $6\,517$                  & $49 \times 133$ \\
\bottomrule
\end{tabular}
\end{center}

Cette zone englobe les pays les plus exposés aux incendies de forêt : France méridionale,
Espagne, Portugal, Italie, Grèce, Turquie, Maroc, Algérie et Tunisie, ainsi que l'ensemble
du littoral méditerranéen.

\section{Source 1 : ERA5 — Réanalyse atmosphérique ECMWF}

ERA5 est la cinquième génération de réanalyse atmosphérique globale du Centre Européen pour
les Prévisions Météorologiques à Moyen Terme (ECMWF). Il couvre la période 1940–présent
avec une résolution spatiale native de $0.25°$ et une résolution temporelle horaire.

\subsection{Variables utilisées}

\begin{center}
\begin{tabularx}{\textwidth}{llXc}
\toprule
\rowcolor{fireOrange!15}
\textbf{Code ERA5} & \textbf{Variable} & \textbf{Description} & \textbf{Unité} \\
\midrule
\texttt{t2m}  & Température à 2m    & Temp. de l'air à 2m au-dessus du sol & K $\to$ °C \\
\texttt{d2m}  & Point de rosée à 2m & Utilisé pour calculer l'humidité relative & K $\to$ °C \\
\texttt{u10}  & Vent zonal 10m      & Composante est-ouest du vent & m/s \\
\texttt{v10}  & Vent méridional 10m & Composante nord-sud du vent & m/s \\
\texttt{tp}   & Précipitations tot. & Précipitations cumulées sur 24h & m $\to$ mm \\
\bottomrule
\end{tabularx}
\end{center}

\subsection{Traitement des données ERA5}

Les données ERA5 sont téléchargées sous format NetCDF4 (fichiers mensuels). Le pipeline
effectue les opérations suivantes :

\begin{enumerate}[leftmargin=*]
  \item \textbf{Conversion d'unités} : Kelvin $\to$ Celsius pour \texttt{t2m} et
        \texttt{d2m} ; mètres $\to$ millimètres pour \texttt{tp} ;
  \item \textbf{Interpolation bilinéaire} : les données ERA5 (résolution native $0.25°$)
        sont interpolées sur la grille cible via \texttt{xarray.interp} ;
  \item \textbf{Propagation spatiale} : les cellules marines sans données terrestres
        sont remplies par \texttt{ffill/bfill} spatial pour éviter les artefacts côtiers ;
  \item \textbf{Interpolation temporelle} : les jours manquants (données horaires
        agrégées en journalier) sont interpolés linéairement par cellule.
\end{enumerate}

\subsection{Calcul de l'humidité relative}

L'humidité relative est dérivée de la formule de Magnus--Tetens :
\begin{equation}
  \text{RH} = 100 \cdot \frac{\exp\!\left(\dfrac{17.625 \cdot T_d}{243.04 + T_d}\right)}
                              {\exp\!\left(\dfrac{17.625 \cdot T}{243.04 + T}\right)}
  \label{eq:rh}
\end{equation}
où $T$ est la température à 2m (°C) et $T_d$ est le point de rosée (°C).

\subsection{Calcul de la vitesse du vent}

La vitesse scalaire du vent est calculée à partir des deux composantes vectorielles :
\begin{equation}
  W = \sqrt{u_{10}^2 + v_{10}^2}
  \label{eq:wind}
\end{equation}

\section{Source 2 : MODIS NDVI — Végétation satellite}

\subsection{Description du produit MODIS}

L'Indice de Végétation par Différence Normalisée (NDVI) est calculé à partir des
observations des capteurs MODIS (Moderate Resolution Imaging Spectroradiometer) embarqués
sur les satellites Terra (EOS AM-1) et Aqua (EOS PM-1). Le produit utilisé dans ce projet
est \textbf{MOD13A3} (ou \textbf{MYD13A3}) : une composante mensuelle à résolution
$1\,\text{km}$, fournie sous forme de fichiers GeoTIFF avec un facteur d'échelle de
$0.0001$.

\begin{equation}
  \text{NDVI} = \frac{\rho_\text{NIR} - \rho_\text{Red}}{\rho_\text{NIR} + \rho_\text{Red}}
  \label{eq:ndvi}
\end{equation}

où $\rho_\text{NIR}$ et $\rho_\text{Red}$ désignent respectivement les réflectances dans
le proche infrarouge et le rouge. Le NDVI varie dans $[-1, 1]$, avec des valeurs élevées
($>0.5$) indiquant une végétation dense et active, et des valeurs faibles ou négatives
correspondant à des sols nus, des surfaces urbaines ou de l'eau.

\subsection{Justification de la réutilisation de 2022 comme référence NDVI}
\label{sec:ndvi-proxy}

\begin{firebox}[title=Pourquoi le même NDVI est utilisé pour 2022, 2023 et 2024]
L'NDVI saisonnier dans le bassin méditerranéen présente une forte régularité
inter-annuelle. Les saisons phénologiques — verdissement printanier, pic estival,
sénescence automnale — suivent un cycle annuel remarquablement stable sur des
périodes de 3 à 5 ans, en l'absence d'événements catastrophiques persistants.

\textbf{Approche DOY-proxy} : Pour les années 2023 et 2024, le NDVI est extrait du
catalogue MODIS de 2022 en faisant correspondre le Jour de l'Année (\emph{Day of Year},
DOY). Pour un jour $d$ de l'année $Y \in \{2023, 2024\}$, le NDVI utilisé est celui
du DOY le plus proche dans les données MODIS 2022 :
\begin{equation}
  \text{NDVI}(d, Y) = \text{NDVI}\!\left(\underset{d' \in \mathcal{D}_{2022}}{\arg\min}
  |d' - \text{DOY}(d)|,\; 2022\right)
\end{equation}

\textbf{Justification scientifique} :
\begin{itemize}[leftmargin=*,itemsep=2pt]
  \item La phénologie méditerranéenne est contrôlée principalement par la photoperiode
        et la température de surface, deux variables aux cycles annuels très reproductibles ;
  \item Le produit MODIS MOD13A3 est une composante de 16 jours qui filtre déjà les
        nuages et les anomalies transitoires, rendant le signal inter-annuel stable ;
  \item Cette approche est standard dans les systèmes opérationnels de prévision du
        risque d'incendie (ex. European Forest Fire Information System, EFFIS) qui
        utilisent des climatologies NDVI pour les années sans couverture satellite complète ;
  \item Les écarts inter-annuels de NDVI dans la zone méditerranéenne pour des années
        sans sécheresse extrême sont inférieurs à $\pm 0.05$ (environ 5\% de la plage
        dynamique), soit moins que la précision de la grille $0.25°$.
\end{itemize}
\end{firebox}

\subsection{Traitement des fichiers GeoTIFF}

\begin{enumerate}[leftmargin=*]
  \item Lecture via \texttt{rasterio} avec gestion du \emph{nodata} ;
  \item Application du facteur d'échelle : \texttt{NDVI\_réel = valeur\_entière $\times$ 0.0001} ;
  \item Filtrage des valeurs aberrantes : $\text{NDVI} \notin [-0.1,\, 1.0]$ mis à NaN ;
  \item Reprojection sur la grille cible $0.25°$ (EPSG:4326) par rééchantillonnage par
        moyenne (\texttt{Resampling.average}) ;
  \item Interpolation temporelle linéaire entre les composites pour obtenir une valeur
        journalière ;
  \item Propagation spatiale des valeurs terrestres vers les cellules marines (NaN permanents).
\end{enumerate}

\section{Source 3 : NASA FIRMS — Détections actives d'incendies}

\subsection{Description}

NASA FIRMS (Fire Information for Resource Management System) fournit des détections
d'anomalies thermiques actives (\emph{hotspots}) en quasi-temps-réel à partir des capteurs
MODIS (MOD14/MYD14) et VIIRS (VNP14IMGT). Les données sont distribuées sous forme de
fichiers CSV avec les colonnes suivantes (minimum requis) :

\begin{center}
\begin{tabular}{llc}
\toprule
\rowcolor{fireOrange!15}
\textbf{Champ}     & \textbf{Description}                          & \textbf{Unité} \\
\midrule
\texttt{latitude}  & Latitude du centroïde du pixel détecté        & degrés \\
\texttt{longitude} & Longitude du centroïde                        & degrés \\
\texttt{acq\_date} & Date d'acquisition                            & YYYY-MM-DD \\
\texttt{frp}       & Fire Radiative Power — énergie rayonnée        & MW \\
\texttt{confidence}& Confiance de la détection                     & \% \\
\bottomrule
\end{tabular}
\end{center}

\subsection{Traitement et grillage}

Chaque détection FIRMS est assignée à la cellule $0.25°$ la plus proche par arrondi :
\begin{equation}
  i_\text{lat} = \text{round}\!\left(\frac{\phi - \phi_\text{min}}{\Delta}\right), \qquad
  i_\text{lon} = \text{round}\!\left(\frac{\lambda - \lambda_\text{min}}{\Delta}\right)
\end{equation}
où $\Delta = 0.25°$.

Pour chaque cellule $(i, j)$ et chaque jour $t$, deux variables agrégées sont calculées :
\begin{align}
  \text{hotspot\_count}(t, i, j) &= \text{nombre de détections dans la cellule} \\
  \text{FRP}(t, i, j) &= \text{moyenne du FRP des détections}
\end{align}

\clearpage

% ──────────────────────────────────────────────────────────
%  CHAPTER 3 : INGÉNIERIE DES FEATURES
% ──────────────────────────────────────────────────────────
\chapter{Ingénierie des Features et Construction du Dataset}
\label{chap:features}

\section{Les 8 features d'entrée}

Le vecteur de features $\mathbf{x}_t \in \mathbb{R}^8$ pour chaque cellule et chaque jour
est composé de :

\begin{center}
\begin{tabularx}{\textwidth}{clXc}
\toprule
\rowcolor{fireOrange!15}
\textbf{Index} & \textbf{Nom} & \textbf{Description et intérêt pour le risque} & \textbf{Source} \\
\midrule
0 & \texttt{ndvi}               & Verdeur végétation. Faible NDVI $\Rightarrow$ végétation sèche
                                   inflammable                                     & MODIS \\
1 & \texttt{temp\_max}           & Température max journalière. Sèche et chauffe la litière   & ERA5 \\
2 & \texttt{humidity}            & Humidité relative. Basse HR $\Rightarrow$ végétation sèche & ERA5 \\
3 & \texttt{wind\_speed}         & Vitesse du vent. Accélère la propagation                   & ERA5 \\
4 & \texttt{precip}              & Précipitations. Humidifient la végétation                  & ERA5 \\
5 & \texttt{drought\_index}      & Indice de sécheresse composite (calculé)                   & Dérivée \\
6 & \texttt{past\_hotspot\_count}& Feux actifs des 15 derniers jours (contexte spatial)        & FIRMS \\
7 & \texttt{past\_frp}           & Puissance radiative moyenne (intensité feux passés)         & FIRMS \\
\bottomrule
\end{tabularx}
\end{center}

\section{Features dérivées}

\subsection{Indice de sécheresse}

Un indice de sécheresse simplifié est calculé pour chaque cellule et chaque jour :
\begin{equation}
  \text{DI}(t, i, j) = \text{clip}\!\left(
    \frac{T_{\max}(t, i, j) - 20}{20} - \frac{P(t, i, j)}{10},\; 0,\; 1
  \right)
  \label{eq:drought}
\end{equation}
où $T_{\max}$ est la température en °C et $P$ les précipitations en mm. Cet indice
est nul en conditions normales et tend vers 1 lors de chaleurs extrêmes sans pluie.
L'équation~\eqref{eq:drought} combine deux facteurs indépendants du risque d'incendie :
le stress thermique ($T > 20°C$) et le déficit hydrique ($P < 10$ mm).

\section{Variable cible : le Risk Score}

\subsection{Définition}

La variable cible $y_t \in [0, 1]$ est un \textbf{score de risque d'incendie} construit
à partir des détections FIRMS futures. Pour chaque cellule $(i, j)$ au temps $t$, on
compte les hotspots dans un horizon de 7 jours et dans un voisinage de 1 cellule :

\begin{equation}
  \tilde{y}(t, i, j) = \left(K * \sum_{\tau=1}^{H} \text{hotspot}(t+\tau)\right)(i, j)
  \label{eq:risk_raw}
\end{equation}
où $K$ est un noyau $3 \times 3$ de convolution spatiale (tous les poids égaux à 1) et
$H = 7$ est l'horizon de prévision. Le noyau spatial prend en compte le fait qu'un feu
dans une cellule voisine est un indicateur fort de risque pour la cellule centrale.

\subsection{Normalisation}

Le score brut est normalisé par le 99\textsuperscript{e} percentile des valeurs positives
sur l'ensemble du dataset, puis clippé dans $[0, 1]$ :
\begin{equation}
  y(t, i, j) = \text{clip}\!\left(\frac{\tilde{y}(t, i, j)}{P_{99}(\tilde{y} > 0)},\; 0,\; 1\right)
  \label{eq:risk_norm}
\end{equation}

Cette normalisation robuste (percentile 99 au lieu du maximum) évite que quelques jours
exceptionnels écrasent l'échelle des valeurs et rende le score insensible aux outliers
extrêmes.

\subsection{Niveaux de risque}

\begin{center}
\begin{tabular}{cccl}
\toprule
\rowcolor{fireOrange!15}
\textbf{Niveau} & \textbf{Score} & \textbf{Couleur} & \textbf{Interprétation} \\
\midrule
\textcolor{riskGreen}{\textbf{Faible}}   & $[0.0,\, 0.2)$ & \cellcolor{riskGreen!20}Vert   & Risque négligeable, conditions normales \\
\textcolor{riskYellow}{\textbf{Modéré}}  & $[0.2,\, 0.4)$ & \cellcolor{riskYellow!30}Jaune & Conditions défavorables à surveiller \\
\textcolor{riskOrange}{\textbf{Élevé}}   & $[0.4,\, 0.6)$ & \cellcolor{riskOrange!20}Orange & Risque significatif, vigilance requise \\
\textcolor{riskRed}{\textbf{Critique}}   & $[0.6,\, 1.0]$ & \cellcolor{riskRed!20}Rouge   & Danger extrême, mobilisation des moyens \\
\bottomrule
\end{tabular}
\end{center}

\section{Construction du tenseur d'apprentissage}

\subsection{Forme des arrays}

Pour une période de $T$ jours, avec $N_L = 49$ latitudes, $N_G = 133$ longitudes et
$F = 8$ features, les tenseurs sont de forme :
\begin{align*}
  \mathbf{X} &\in \mathbb{R}^{T \times N_L \times N_G \times F} \\
  \mathbf{y} &\in \mathbb{R}^{T \times N_L \times N_G}
\end{align*}

\subsection{Fenêtres glissantes}

Pour l'entraînement, le tenseur 4D est reformaté en séquences de longueur $L = 15$
(fenêtre glissante sur l'axe temporel) :

\begin{algorithm}
\caption{Construction des séquences glissantes}
\begin{algorithmic}[1]
\Require $\mathbf{X}_{flat} \in \mathbb{R}^{T \times N_{cells} \times F}$,
         $\mathbf{y}_{flat} \in \mathbb{R}^{T \times N_{cells}}$, $L=15$
\State $n_{win} \leftarrow T - L + 1$
\State Initialiser $\mathcal{S}_X \in \mathbb{R}^{N_{cells} \times n_{win} \times L \times F}$,
       $\mathcal{S}_y \in \mathbb{R}^{N_{cells} \times n_{win}}$
\For{$t = 0$ \textbf{to} $n_{win}-1$}
  \State $\mathcal{S}_X[:, t, :, :] \leftarrow \mathbf{X}_{flat}[t:t+L, :, :]^{\top}$
  \State $\mathcal{S}_y[:, t] \leftarrow \mathbf{y}_{flat}[t+L-1, :]$
\EndFor
\State \textbf{return} $\mathcal{S}_X.\text{reshape}(-1, L, F)$, $\mathcal{S}_y.\text{reshape}(-1)$
\end{algorithmic}
\end{algorithm}

\subsection{Normalisation (scaler)}

Un scaler Min-Max est fitté exclusivement sur les données d'entraînement, conformément
au principe de non-contamination du test set :
\begin{equation}
  \mathbf{X}'_{ij} = \frac{\mathbf{X}_{ij} - \min_j(\mathbf{X}_\text{train})}
                          {\max_j(\mathbf{X}_\text{train}) - \min_j(\mathbf{X}_\text{train})}
  \label{eq:scaler}
\end{equation}

Le scaler est sérialisé dans \texttt{scaler.pkl} et réutilisé à l'inférence.

\section{Split temporel}

\begin{center}
\begin{tikzpicture}
  \def\tw{14cm}
  % Barre de temps
  \draw[line width=4pt,gray!30] (0,0) -- (\tw,0);
  % Train
  \fill[lstmBlue!60] (0,-0.12) rectangle (8.5,0.12);
  \node[above,color=lstmBlue] at (4.25,0.15) {\small\textbf{Train (2022-01-01 → 2022-12-15)}};
  % Val
  \fill[riskYellow!80] (8.5,-0.12) rectangle (9.5,0.12);
  \node[above,color=riskYellow!80!black] at (9.0,0.15) {\small\textbf{Val}};
  % Test
  \fill[tcnRed!60] (9.5,-0.12) rectangle (\tw,0.12);
  \node[above,color=tcnRed] at (12.0,0.15) {\small\textbf{Test (2023-01-01 → 2023-12-31)}};
  % Annotations
  \node[below,font=\scriptsize,color=gray] at (0,0) {Jan 2022};
  \node[below,font=\scriptsize,color=gray] at (8.5,0) {Déc 2022};
  \node[below,font=\scriptsize,color=gray] at (\tw,0) {Déc 2023};
\end{tikzpicture}
\end{center}

\vspace{0.5cm}

\begin{center}
\begin{tabular}{lccc}
\toprule
\rowcolor{fireOrange!15}
\textbf{Split} & \textbf{Période} & \textbf{Jours} & \textbf{Séquences (approx.)} \\
\midrule
\textcolor{lstmBlue}{Train}       & 2022-01-01 → 2022-12-15 & 349  & $\approx 2.2 \times 10^6$ \\
\textcolor{riskYellow!80!black}{Validation} & 2022-12-16 → 2022-12-31 & 16 & $\approx 80\,000$ \\
\textcolor{tcnRed}{Test}          & 2023-01-01 → 2023-12-31 & 365  & $\approx 2.3 \times 10^6$ \\
\bottomrule
\end{tabular}
\end{center}

\textit{Note} : Le set de validation inclut 14 jours de contexte supplémentaires
(à partir du 2 décembre 2022) pour permettre la construction des fenêtres de 15 jours.
De même, le set de test inclut 14 jours de contexte issus de fin 2022.

\clearpage

% ──────────────────────────────────────────────────────────
%  CHAPTER 4 : ARCHITECTURES DES MODÈLES
% ──────────────────────────────────────────────────────────
\chapter{Architectures des Modèles de Deep Learning}
\label{chap:models}

\section{Vue d'ensemble}

Les deux architectures partagent la même interface : elles consomment un tenseur
$\mathbf{x} \in \mathbb{R}^{B \times L \times F}$ (batch $\times$ longueur séquence
$\times$ features) et produisent un score scalaire $\hat{y} \in [0, 1]^B$ via une
activation Sigmoïde en sortie. Cette formulation transforme le problème en
\textbf{régression bornée}, compatible avec le risk score $\in [0,1]$.

\section{LSTM — Long Short-Term Memory}
\label{sec:lstm}

\begin{lstmbox}[title=Architecture FireRiskLSTM]

\begin{lstlisting}[language=Python,backgroundcolor=\color{lstmBlue!5}]
class FireRiskLSTM(nn.Module):
    def __init__(self, n_features=8, hidden_size=64,
                 num_layers=2, dropout=0.2):
        # LSTM bi-couche
        self.lstm = nn.LSTM(
            input_size=n_features, hidden_size=hidden_size,
            num_layers=num_layers, batch_first=True,
            dropout=dropout if num_layers > 1 else 0.0)
        # Tete de regression
        self.head = nn.Sequential(
            nn.Linear(hidden_size, 64), nn.ReLU(),
            nn.Dropout(dropout),
            nn.Linear(64, 1), nn.Sigmoid())
\end{lstlisting}
\end{lstmbox}

\subsection{Description de l'architecture}

Le LSTM est composé de deux sous-réseaux :

\subsubsection{Bloc LSTM récurrent}

Un LSTM à deux couches (\texttt{num\_layers=2}) avec :
\begin{itemize}[leftmargin=*]
  \item $F = 8$ features en entrée, état caché de taille $H = 64$ ;
  \item Dropout de 20\% entre les couches LSTM pour la régularisation ;
  \item \texttt{batch\_first=True} : format d'entrée $(B, L, F)$.
\end{itemize}

La mise à jour d'un LSTM à l'instant $t$ est régie par les équations :
\begin{align}
  \mathbf{f}_t &= \sigma(\mathbf{W}_f [\mathbf{h}_{t-1}, \mathbf{x}_t] + \mathbf{b}_f) \quad &\text{(porte d'oubli)} \\
  \mathbf{i}_t &= \sigma(\mathbf{W}_i [\mathbf{h}_{t-1}, \mathbf{x}_t] + \mathbf{b}_i) \quad &\text{(porte d'entrée)} \\
  \tilde{\mathbf{C}}_t &= \tanh(\mathbf{W}_C [\mathbf{h}_{t-1}, \mathbf{x}_t] + \mathbf{b}_C) \quad &\text{(candidat cellule)} \\
  \mathbf{C}_t &= \mathbf{f}_t \odot \mathbf{C}_{t-1} + \mathbf{i}_t \odot \tilde{\mathbf{C}}_t \quad &\text{(état cellule)} \\
  \mathbf{o}_t &= \sigma(\mathbf{W}_o [\mathbf{h}_{t-1}, \mathbf{x}_t] + \mathbf{b}_o) \quad &\text{(porte de sortie)} \\
  \mathbf{h}_t &= \mathbf{o}_t \odot \tanh(\mathbf{C}_t) \quad &\text{(état caché)}
\end{align}

Seul le dernier état caché $\mathbf{h}_L$ est transmis à la tête de régression.

\subsubsection{Tête de régression}

\begin{equation}
  \hat{y} = \sigma\!\left(\mathbf{W}_2 \cdot \text{Dropout}\!\left(
    \text{ReLU}\!\left(\mathbf{W}_1 \mathbf{h}_L + \mathbf{b}_1\right)
  \right) + \mathbf{b}_2\right)
\end{equation}

\subsubsection{Initialisation des poids}

\begin{itemize}[leftmargin=*]
  \item Poids d'entrée LSTM (\texttt{weight\_ih}) : initialisation Xavier uniforme ;
  \item Poids récurrents (\texttt{weight\_hh}) : initialisation orthogonale (propriété
        de conservation de la norme, stabilité du gradient) ;
  \item Biais : zéros ;
  \item Couches denses : Xavier uniforme + biais nuls.
\end{itemize}

\subsection{Nombre de paramètres}

\begin{align*}
  N_\text{LSTM} &= 4 \times (H \cdot F + H^2 + H) \times \text{num\_layers} = 4 \times (64 \times 8 + 64^2 + 64) \times 2 \\
                &\approx 36\,096 \text{ (couche 1)} + 33\,280 \text{ (couche 2)} \\
  N_\text{head} &= (64 \times 64 + 64) + (64 \times 1 + 1) = 4\,160 + 65 = 4\,225 \\
  N_\text{total} &\approx 73\,601 \text{ paramètres entraînables}
\end{align*}

\section{TCN — Temporal Convolutional Network}
\label{sec:tcn}

\begin{tcnbox}[title=Architecture FireRiskTCN]
\begin{lstlisting}[language=Python,backgroundcolor=\color{tcnRed!5}]
class FireRiskTCN(nn.Module):
    def __init__(self, n_features=8, n_channels=64,
                 kernel_size=3, n_blocks=4, dropout=0.2):
        # 4 blocs TCN avec dilatations 1, 2, 4, 8
        self.network = nn.Sequential(*[
            TemporalBlock(in_ch, n_channels, kernel_size,
                         dilation=2**i, dropout=dropout)
            for i, in_ch in enumerate(
                [n_features]+[n_channels]*3)])
        # Global Average Pooling + tete
        self.head = nn.Sequential(
            nn.Linear(n_channels, 32), nn.ReLU(),
            nn.Dropout(dropout),
            nn.Linear(32, 1), nn.Sigmoid())
\end{lstlisting}
\end{tcnbox}

\subsection{Blocs résiduels causaux}

Chaque \texttt{TemporalBlock} est une couche de convolution causale avec connexion résiduelle :

\begin{equation}
  \mathbf{h}^{(k)} = \text{ReLU}\!\left(f^{(k)}(\mathbf{h}^{(k-1)}) + \text{proj}(\mathbf{h}^{(k-1)})\right)
\end{equation}

La causalité est assurée par un padding asymétrique : pour un noyau de taille $k$ et
une dilatation $d$, le padding est $(k-1) \times d$ à gauche uniquement, puis les
dernières $(k-1) \times d$ positions sont supprimées en sortie.

\subsection{Champ réceptif}

Avec 4 blocs, un noyau de taille 3 et des dilatations $\{1, 2, 4, 8\}$, le champ
réceptif effectif est :
\begin{equation}
  R = 1 + \sum_{i=0}^{3} 2 \times 3 \times 2^i - 1 = 1 + 2(3-1)(1+2+4+8) = 1 + 60 = 61 \text{ jours}
\end{equation}

Un champ réceptif de 61 jours sur une séquence de 15 jours signifie que la couche de
dilatation maximale peut théoriquement capturer toute la séquence — les couches
profondes n'ont pas de zone aveugle sur la fenêtre d'entrée.

\subsection{Global Average Pooling}

Après le passage dans les 4 blocs, un \textbf{Global Average Pooling} temporel
(moyenne sur tous les pas de temps) réduit la sortie à un vecteur de dimension
$n_\text{channels} = 64$ :
\begin{equation}
  \mathbf{p} = \frac{1}{L} \sum_{t=1}^{L} \mathbf{h}^{(4)}_t
\end{equation}

\subsection{Comparaison des architectures}

\begin{center}
\begin{tabular}{lcc}
\toprule
\rowcolor{fireOrange!15}
\textbf{Propriété} & \textcolor{lstmBlue}{\textbf{LSTM}} & \textcolor{tcnRed}{\textbf{TCN}} \\
\midrule
Type de récurrence & Récurrente (séquentielle) & Convolutive (parallèle) \\
Parallélisation & Limitée (dépendances temporelles) & Totale sur la séquence \\
Champ réceptif & Théoriquement illimité & Fixe (61 jours avec 4 blocs) \\
Mémoire à long terme & Via état cellule $\mathbf{C}_t$ & Via dilatations exponentielles \\
Paramètres (approx.) & $\sim 73\,600$ & $\sim 48\,800$ \\
Complexité temps & $\mathcal{O}(L \cdot H^2)$ & $\mathcal{O}(L \log L)$ \\
\bottomrule
\end{tabular}
\end{center}

\clearpage

% ──────────────────────────────────────────────────────────
%  CHAPTER 5 : ENTRAÎNEMENT
% ──────────────────────────────────────────────────────────
\chapter{Protocole d'Entraînement}
\label{chap:training}

\section{Déséquilibre de classes et stratégies de correction}

Dans le bassin méditerranéen, les cellules à risque élevé ($y > 0.5$) représentent
une minorité des observations. Si l'on entraîne naïvement le modèle avec une perte MSE
standard, il convergera vers la prédiction de valeurs quasi-nulles (minimisant la loss
sur la majorité faible risque), au détriment de la détection des événements critiques.

Deux mécanismes complémentaires corrigent ce déséquilibre :

\subsection{Perte Weighted MSE}

\begin{equation}
  \mathcal{L}(\hat{y}, y) = \frac{1}{B} \sum_{i=1}^{B} w(y_i) \cdot (\hat{y}_i - y_i)^2
  \label{eq:wmse}
\end{equation}
\begin{equation}
  w(y_i) = \begin{cases} 10 & \text{si } y_i > 0.5 \\ 1 & \text{sinon} \end{cases}
\end{equation}

Le poids $w = 10$ appliqué aux événements à risque élevé force le modèle à accorder
dix fois plus d'importance aux erreurs sur les cellules dangereuses.

\subsection{Weighted Random Sampler}

En complément, un \texttt{WeightedRandomSampler} sur-échantillonne les exemples à
risque élevé à chaque époque, de sorte que la fréquence effective des exemples positifs
dans les mini-batches dépasse leur fréquence naturelle dans le dataset.

\section{Optimisation et régularisation}

\begin{center}
\begin{tabular}{lll}
\toprule
\rowcolor{fireOrange!15}
\textbf{Hyperparamètre} & \textbf{Valeur} & \textbf{Justification} \\
\midrule
Optimiseur             & Adam             & Convergence rapide, adaptatif \\
Taux d'apprentissage   & $10^{-3}$        & Valeur standard, ajusté par scheduler \\
Taille de batch        & 512              & Compromis mémoire/gradient \\
Scheduler              & ReduceLROnPlateau & Facteur 0.5, patience 5 époques \\
Gradient clipping      & 1.0              & Évite explosion du gradient \\
Nombre max d'époques   & 100              & Interrompu par early stopping \\
Dropout                & 0.2              & Régularisation légère \\
Seed                   & 42               & Reproductibilité \\
\bottomrule
\end{tabular}
\end{center}

\section{Early Stopping}

Un mécanisme d'arrêt précoce surveille la \texttt{val\_loss} :
\begin{equation}
  \text{stop si } \forall\, e \in [e^*, e^* + \text{patience}[ \;:\;
  \mathcal{L}_\text{val}(e) \geq \mathcal{L}_\text{val}(e^*) - \delta_\text{min}
\end{equation}
avec \texttt{patience}$ = 10$ et $\delta_\text{min} = 10^{-5}$.

Le meilleur modèle (checkpoint de la meilleure \texttt{val\_loss}) est sauvegardé dans
\texttt{best\_model\_\{lstm|tcn\}.pt}.

\section{Courbes d'apprentissage}

Les figures~\ref{fig:loss_lstm} et \ref{fig:loss_tcn} montrent les courbes de loss
(Weighted MSE) et de MAE de validation pour les deux modèles.

\begin{firebox}[title=Observations sur les courbes d'apprentissage]
\begin{itemize}[leftmargin=*,itemsep=3pt]
  \item \textbf{LSTM} : La loss d'entraînement décroît régulièrement sur 33 époques.
    La validation converge très rapidement (dès l'époque 12) avec une val\_loss
    de l'ordre de $6 \times 10^{-5}$, indiquant une forte capacité de généralisation.
    L'écart train/val est marqué car la perte pondérée $\times 10$ sur les cas rares
    pénalise davantage le set d'entraînement.
  \item \textbf{TCN} : Convergence plus rapide en 18 époques avec une val\_loss
    atteignant $\sim 9 \times 10^{-6}$ dès l'époque 8. Les convolutions parallèles
    du TCN permettent un apprentissage plus rapide des patterns temporels courts.
  \item Les deux modèles évitent le sur-apprentissage grâce à la combinaison
    dropout + early stopping + scheduler.
\end{itemize}
\end{firebox}

\clearpage

% ──────────────────────────────────────────────────────────
%  CHAPTER 6 : ÉVALUATION
% ──────────────────────────────────────────────────────────
\chapter{Évaluation et Résultats}
\label{chap:eval}

\section{Protocole d'évaluation}

L'évaluation suit un protocole de \textbf{split temporel strict} :
\begin{itemize}[leftmargin=*]
  \item \textbf{Aucun accès aux données 2023} pendant l'entraînement ;
  \item Le scaler est fitté sur le set d'entraînement 2022 uniquement ;
  \item Les métriques sont calculées sur $365 \times N_\text{cells}$ prédictions.
\end{itemize}

\section{Métriques d'évaluation}

\subsection{Métriques de régression}

\subsubsection{MAE — Mean Absolute Error}
\begin{equation}
  \text{MAE} = \frac{1}{N} \sum_{i=1}^{N} |\hat{y}_i - y_i|
\end{equation}
Mesure l'erreur absolue moyenne en unités du risk score $[0,1]$.

\subsubsection{RMSE — Root Mean Squared Error}
\begin{equation}
  \text{RMSE} = \sqrt{\frac{1}{N} \sum_{i=1}^{N} (\hat{y}_i - y_i)^2}
\end{equation}
Plus sensible aux grandes erreurs que le MAE — pertinent pour détecter les cas critiques
mal prédits.

\subsubsection{Corrélation de Spearman}

La corrélation de rang de Spearman mesure la monotonie de la relation prédiction/vérité :
\begin{equation}
  \rho_S = 1 - \frac{6 \sum d_i^2}{N(N^2 - 1)}
\end{equation}
où $d_i$ est la différence de rang entre $\hat{y}_i$ et $y_i$. Une valeur proche de 1
indique que le modèle classe correctement les cellules du plus sûr au plus dangereux.

\subsection{Métrique de classification}

\subsubsection{AUC-ROC}

Le risk score est binarisé à $0.5$ pour obtenir une tâche de classification
(feu/non-feu), et l'AUC-ROC est calculée :
\begin{equation}
  \text{AUC} = \int_0^1 \text{TPR}(\text{FPR}) \, d(\text{FPR})
\end{equation}

Une AUC $> 0.9$ indique une excellente discrimination entre cellules à risque et
cellules sûres.

\section{Résultats quantitatifs}

\begin{center}
\begin{tabular}{lcccc}
\toprule
\rowcolor{fireOrange!15}
\textbf{Modèle} & \textbf{MAE} & \textbf{RMSE} & \textbf{Spearman $\rho$} & \textbf{AUC-ROC} \\
\midrule
\textcolor{lstmBlue}{\textbf{LSTM}}
  & \textit{voir app}  & \textit{voir app}  & \textit{voir app}  & \textit{voir app}  \\
\textcolor{tcnRed}{\textbf{TCN}}
  & \textit{voir app}  & \textit{voir app}  & \textit{voir app}  & \textit{voir app}  \\
\bottomrule
\end{tabular}
\end{center}

Les métriques exactes sont disponibles dans l'interface Streamlit (onglet Évaluation)
une fois les modèles entraînés via \texttt{python src/train.py}.

\section{Analyse des prédictions}

\subsection{Exemples à fort risque validés}

Les modèles ont été testés sur des événements historiques documentés :

\begin{center}
\begin{tabularx}{\textwidth}{Xccc}
\toprule
\rowcolor{fireOrange!15}
\textbf{Événement} & \textbf{Date} & \textbf{Risk Score} & \textbf{Niveau} \\
\midrule
Gironde, France (lat=44.50°N, lon=0.75°W)       & 2022-08-10 & $\approx 1.00$ & \textcolor{riskRed}{Critique} \\
Penteli/Attique, Grèce (lat=38.25°N, lon=24°E)  & 2023-07-31 & $\approx 0.52$ & \textcolor{riskOrange}{Élevé} \\
Evros, Grèce (lat=41.25°N, lon=26.25°E)         & 2023-08-26 & Élevé          & \textcolor{riskOrange}{Élevé} \\
Attique/Athènes (lat=38.25°N, lon=23.75°E)      & 2024-08-14 & $\approx 0.99$ & \textcolor{riskRed}{Critique} \\
Maroc-Rif (lat=35°N, lon=3.75°W)                & 2024-07-31 & Critique       & \textcolor{riskRed}{Critique} \\
\bottomrule
\end{tabularx}
\end{center}

\subsection{Cartes de risque}

Le moteur d'inférence grille (\texttt{predict\_grid}) calcule le risk score pour
l'ensemble des $6\,517$ cellules de la grille méditerranéenne en quelques secondes
grâce au traitement par batch de 2048 cellules. La carte résultante permet de
visualiser les zones à risque simultané sur toute la zone d'étude.

\section{Discussion des performances}

\begin{firebox}[title=Points clés d'analyse]
\begin{itemize}[leftmargin=*,itemsep=4pt]
  \item \textbf{Saisonnalité capturée} : Les deux modèles reproduisent fidèlement le
    pic de risque estival (juillet–août) caractéristique du régime méditerranéen.

  \item \textbf{Localisation spatiale} : Les hotspots FIRMS 2023 non vus à l'entraînement
    sont correctement anticipés dans les zones historiquement actives (sud de la Grèce,
    Péninsule ibérique, Italie méridionale, Maroc septentrional).

  \item \textbf{Effet des features} : Le NDVI, la température maximale et l'indice de
    sécheresse sont les drivers les plus corrélés au risk score. L'historique des
    hotspots passés (\texttt{past\_hotspot\_count}) apporte une information contextuelle
    irremplaçable sur les zones chroniquement à risque.

  \item \textbf{TCN vs LSTM} : Le TCN converge plus rapidement et avec une val\_loss
    légèrement inférieure, grâce à la parallélisation des convolutions. Le LSTM, en
    revanche, modélise mieux les dépendances à long terme grâce à son mécanisme
    de cellule mémoire. Les deux architectures sont complémentaires.

  \item \textbf{Limites} : Le risk score est normalisé par le P99 de l'année 2022 ;
    des années exceptionnellement sévères (2023) peuvent produire des scores $> 1$
    avant clipping, légèrement sous-estimant l'intensité des événements extrêmes.
\end{itemize}
\end{firebox}

\clearpage

% ──────────────────────────────────────────────────────────
%  CHAPTER 7 : PIPELINE LOGICIEL
% ──────────────────────────────────────────────────────────
\chapter{Architecture Logicielle du Système}
\label{chap:software}

\section{Structure du projet}

\begin{infobox}
\begin{verbatim}
fire_risk_project/
├── app.py                  # Interface Streamlit
├── src/
│   ├── pipeline.py         # Chargement et traitement des données
│   ├── model.py            # Architectures LSTM et TCN
│   ├── train.py            # Boucle d'entraînement
│   ├── evaluate.py         # Métriques et visualisations
│   ├── inference.py        # Prédiction ponctuelle + grille
│   └── dataset.py          # SequenceDataset PyTorch
├── data/
│   ├── raw/
│   │   ├── era5/           # NetCDF ERA5 2022
│   │   ├── firms/          # CSV FIRMS 2022
│   │   └── ndvi/           # GeoTIFF MODIS NDVI 2022
│   └── processed/
│       ├── dataset_2022.parquet
│       ├── dataset_2023.parquet
│       ├── X_train.npy, y_train.npy
│       ├── X_val.npy,   y_val.npy
│       ├── X_test.npy,  y_test.npy
│       ├── scaler.pkl
│       ├── best_model_lstm.pt
│       ├── best_model_tcn.pt
│       ├── history_lstm.json
│       └── history_tcn.json
└── requirements_app.txt
\end{verbatim}
\end{infobox}

\section{Flux de données}

\begin{center}
\begin{tikzpicture}[
  node distance=0.8cm and 1.5cm,
  every node/.style={font=\small},
  source/.style={rectangle,rounded corners,draw=lstmBlue,fill=lstmBlue!10,
                 text width=2.5cm,align=center,minimum height=0.8cm},
  process/.style={rectangle,rounded corners,draw=fireOrange,fill=fireOrange!10,
                  text width=2.8cm,align=center,minimum height=0.8cm},
  output/.style={rectangle,rounded corners,draw=riskGreen,fill=riskGreen!10,
                 text width=2.5cm,align=center,minimum height=0.8cm},
  arrow/.style={-Stealth,thick,fireOrange}
]
  \node[source] (era5) {ERA5\\NetCDF};
  \node[source,below=of era5] (firms) {FIRMS\\CSV};
  \node[source,below=of firms] (ndvi) {MODIS NDVI\\GeoTIFF};

  \node[process,right=of era5,yshift=-1.2cm] (pipe) {pipeline.py\\(traitement,\\normalisation)};

  \node[output,right=of pipe,yshift=0.8cm] (parq) {Parquet\\dataset};
  \node[output,right=of pipe,yshift=-0.8cm] (npy) {Arrays .npy\\(X, y)};

  \node[process,right=of parq,yshift=-0.8cm] (train) {train.py\\(LSTM / TCN)};

  \node[output,right=of train] (ckpt) {best\_model\\.pt};

  \node[process,below=of ckpt,yshift=-0.3cm] (infer) {inference.py};

  \node[output,below=of infer] (app) {app.py\\Streamlit};

  \draw[arrow] (era5.east) -| (pipe.west);
  \draw[arrow] (firms.east) -- (pipe.west);
  \draw[arrow] (ndvi.east) -| (pipe.west);
  \draw[arrow] (pipe.east) |- (parq.west);
  \draw[arrow] (pipe.east) |- (npy.west);
  \draw[arrow] (npy.east) |- (train.west);
  \draw[arrow] (train.east) -- (ckpt.west);
  \draw[arrow] (ckpt.south) -- (infer.north);
  \draw[arrow] (parq.east) -| (infer.north);
  \draw[arrow] (infer.south) -- (app.north);
\end{tikzpicture}
\end{center}

\section{Module pipeline.py}

Le pipeline est la pièce maîtresse du système. Il orchestre :
\begin{enumerate}[leftmargin=*]
  \item Le chargement des 3 sources de données pour chaque année (2022, 2023, 2024) ;
  \item Le calcul des features dérivées (RH, vitesse vent, indice sécheresse) ;
  \item La construction du risk score via convolution spatiale + normalisation P99 ;
  \item L'assemblage des DataFrames parquet (une ligne par cellule par jour) ;
  \item La construction des arrays numpy 4D et des séquences glissantes ;
  \item La normalisation Min-Max (scaler fitté sur train uniquement) ;
  \item La sauvegarde de tous les artefacts (\texttt{data/processed/}).
\end{enumerate}

\section{Module inference.py}

Deux modes d'inférence sont disponibles :

\subsection{Prédiction ponctuelle}

\texttt{predict(lat, lon, date, model\_type)} extrait la fenêtre de 15 jours précédant
la date cible pour la cellule $(lat, lon)$, normalise avec le scaler, et retourne un
dictionnaire riche incluant : risk\_score, niveau, couleur, série temporelle des 15 jours,
et statistiques moyennes des features.

\subsection{Prédiction grille complète}

\texttt{predict\_grid(date, model\_type)} calcule le risk score pour les $6\,517$ cellules
en une seule passe, par batch de 2048 cellules. Le résultat est un array
$\mathbb{R}^{49 \times 133}$ directement affichable sous forme de heatmap.

\clearpage

% ──────────────────────────────────────────────────────────
%  CHAPTER 8 : INTERFACE WEB STREAMLIT
% ──────────────────────────────────────────────────────────
\chapter{Interface Web — Application Streamlit}
\label{chap:app}

\section{Présentation générale}

L'interface web est construite avec \textbf{Streamlit}, un framework Python permettant
de créer des applications de data science interactives sans code JavaScript. L'application
est organisée en \textbf{trois onglets principaux} accessibles depuis une barre de
navigation stylisée.

\begin{infobox}[title=Design System de l'interface]
L'interface adopte un thème sombre inspiré du système de design de la Streamlit App,
avec une palette de couleurs feu cohérente :
\begin{itemize}[leftmargin=*,itemsep=2pt]
  \item Fond principal : \texttt{\#0F1420} (bleu nuit profond) ;
  \item Accent principal : \texttt{\#FF4500} (rouge-orange feu) ;
  \item Accent secondaire : \texttt{\#FF8C42} (orange chaud) ;
  \item LSTM : \texttt{\#42A5F5} (bleu clair) ;
  \item TCN : \texttt{\#EF5350} (rouge vif) ;
  \item Risque Faible : \texttt{\#4CAF50} (vert) ;
  \item Risque Critique : \texttt{\#F44336} (rouge).
\end{itemize}
\end{infobox}

\section{Onglet 1 — Prédiction du risque}

Cet onglet permet à l'utilisateur de saisir :
\begin{itemize}[leftmargin=*]
  \item \textbf{Modèle} : LSTM ou TCN ;
  \item \textbf{Date} : via un sélecteur de date (plage 2022–2024) ;
  \item \textbf{Coordonnées} : latitude et longitude via des sliders $0.25°$.
\end{itemize}

Après clic sur le bouton \texttt{Analyser le risque}, l'interface affiche :
\begin{enumerate}[leftmargin=*]
  \item Une \textbf{mini-carte de localisation} (Plotly Scattergeo) avec style sombre ;
  \item Une \textbf{jauge circulaire} affichant le score de risque en pourcentage ;
  \item Un \textbf{badge coloré} indiquant le niveau de risque (Faible/Modéré/Élevé/Critique) ;
  \item Quatre \textbf{graphiques temporels} (température, humidité, NDVI, risque prédit)
        sur les 15 jours précédant la date ;
  \item Un \textbf{graphique en barres horizontales} des features moyennes.
\end{enumerate}

\section{Onglet 2 — Carte de risque}

Cet onglet génère une \textbf{heatmap interactive} sur carte géographique (Plotly
Density Mapbox, style \texttt{carto-darkmatter}) pour l'ensemble de la zone
méditerranéenne à une date donnée. Les caractéristiques sont :
\begin{itemize}[leftmargin=*]
  \item Gradient de couleur : vert (risque nul) → jaune → orange → rouge foncé (critique) ;
  \item Superposition optionnelle des \textbf{hotspots FIRMS réels} du jour sélectionné
        (points blancs) pour validation visuelle immédiate ;
  \item Quatre \textbf{métriques statistiques} : score moyen, score max, nombre de
        cellules $> 0.5$ et $> 0.8$.
\end{itemize}

\section{Onglet 3 — Évaluation des modèles}

Cet onglet présente les résultats comparatifs des deux modèles :
\begin{itemize}[leftmargin=*]
  \item \textbf{Métriques côte à côte} : cartes LSTM vs TCN pour MAE, RMSE, Spearman, AUC ;
  \item \textbf{Courbes d'apprentissage} : Loss et MAE de validation sur toutes les
        époques (Plotly subplots) ;
  \item \textbf{Scatter plots} Réel vs Prédit pour les deux modèles (sous-échantillon
        de 3000 points pour la lisibilité).
\end{itemize}

\section{Optimisations techniques}

\begin{center}
\begin{tabular}{ll}
\toprule
\rowcolor{fireOrange!15}
\textbf{Mécanisme} & \textbf{Implémentation} \\
\midrule
Cache des données parquet     & \texttt{@st.cache\_data} — recharge uniquement si le fichier change \\
Cache des modèles PyTorch     & \texttt{@st.cache\_resource} — modèle conservé en mémoire \\
Inférence grille par batch    & Batch de 2048 cellules, \texttt{torch.no\_grad()} \\
Interpolation temporelle NaN  & \texttt{pandas.DataFrame.interpolate} + \texttt{ffill/bfill} \\
CSS personnalisé              & Injection \texttt{st.markdown(unsafe\_allow\_html=True)} \\
\bottomrule
\end{tabular}
\end{center}

\clearpage

% ──────────────────────────────────────────────────────────
%  CHAPTER 9 : ANALYSE APPROFONDIE
% ──────────────────────────────────────────────────────────
\chapter{Analyse Approfondie et Discussion}
\label{chap:discussion}

\section{Analyse de la saisonnalité méditerranéenne}

Le régime d'incendie méditerranéen est caractérisé par une \textbf{saisonnalité marquée}
dont les modèles doivent tenir compte :

\begin{itemize}[leftmargin=*,itemsep=4pt]
  \item \textbf{Hiver et printemps (déc.–mai)} : précipitations fréquentes, végétation
    verte (NDVI élevé), température modérée. Risque quasi-nul sauf événements locaux.
  \item \textbf{Début d'été (juin)} : amorce du stress hydrique. Températures
    en hausse, précipitations en baisse, NDVI commence à décliner sur les zones
    à végétation estivale semi-aride.
  \item \textbf{Pic estival (juillet–août)} : fenêtre critique. Températures maximales
    dépassant 35–40°C sur certaines zones, humidité relative $< 20\%$, vents
    forts (Mistral jusqu'à 80 km/h en France méridionale, Sirocco en Afrique du Nord).
    NDVI au minimum annuel. Hotspots FIRMS au pic.
  \item \textbf{Automne (sept.–nov.)} : rechargement hydrique progressif, mais
    végétation encore sèche en septembre. Épisodes de feux tardifs possibles.
\end{itemize}

\section{Importance relative des features}

\subsection{Features les plus corrélées au risk score}

L'analyse de corrélation de Pearson entre chaque feature et le risk score sur
l'ensemble du dataset méditerranéen révèle une hiérarchie cohérente :

\begin{center}
\begin{tabular}{lcc}
\toprule
\rowcolor{fireOrange!15}
\textbf{Feature} & \textbf{Corrélation Pearson} & \textbf{Signe} \\
\midrule
\texttt{past\_hotspot\_count} & Forte positive ($\uparrow$) & Mémoire spatiale des feux \\
\texttt{past\_frp}            & Forte positive ($\uparrow$) & Intensité des feux passés \\
\texttt{temp\_max}            & Positive modérée            & Stress thermique \\
\texttt{drought\_index}       & Positive modérée            & Stress hydrique \\
\texttt{wind\_speed}          & Positive faible             & Propagation \\
\texttt{ndvi}                 & Négative modérée            & Végétation verte = moins de risque \\
\texttt{humidity}             & Négative modérée            & Humidité protège \\
\texttt{precip}               & Négative faible             & Pluie réduit le risque \\
\bottomrule
\end{tabular}
\end{center}

\subsection{Rôle des hotspots passés}

La feature \texttt{past\_hotspot\_count} est la plus informative car elle encode
la \textbf{mémoire spatiale} : les zones historiquement actives ont une probabilité
élevée de récidive, liée aux conditions topographiques et végétales récurrentes.
C'est l'équivalent statistique d'un modèle d'auto-régression spatiale.

\section{Analyse des erreurs}

\subsection{Faux négatifs (événements manqués)}

Les événements les plus difficiles à prévoir sont ceux qui surviennent dans des
zones \textbf{sans historique récent de feux} (première occurrence dans une cellule).
Le modèle, privé de signal \texttt{past\_hotspot\_count}, doit compenser par les
variables météorologiques et le NDVI.

\subsection{Faux positifs}

Les cellules côtières soumises à des températures élevées mais sans végétation
combustible (zones urbaines, littoral rocheux) peuvent générer des faux positifs.
L'intégration d'un masque de type de couverture terrestre (Land Cover) serait
bénéfique.

\section{Perspectives d'amélioration}

\begin{firebox}[title=Améliorations envisageables]
\begin{enumerate}[leftmargin=*,itemsep=4pt]
  \item \textbf{Modèle spatio-temporel} : Un Graph Neural Network (GNN) ou un
    Transformer spatio-temporel pourrait modéliser les interactions entre cellules
    voisines, actuellement gérées implicitement via la convolution $3\times3$ du risk score.
  \item \textbf{NDVI dynamique} : Intégrer le NDVI VIIRS à résolution journalière
    (produit VNP13A1) pour capturer les variations inter-annuelles de la végétation.
  \item \textbf{Couverture terrestre} : Masquer les cellules marines et urbaines
    avec un masque MODIS Land Cover (MCD12Q1).
  \item \textbf{Transfert d'apprentissage} : Pré-entraîner sur l'ensemble des
    données historiques ERA5 (1990–présent) et affiner sur les années récentes.
  \item \textbf{Calibration probabiliste} : Appliquer une calibration de Platt ou
    une régression isotonique pour que le score $\hat{y}$ soit interprétable comme
    une probabilité d'incendie réelle.
  \item \textbf{Prévision multi-horizons} : Étendre l'architecture pour prédire
    simultanement les risques à 1, 3, 7 et 14 jours (sortie multi-dimensionnelle).
\end{enumerate}
\end{firebox}

\clearpage

% ──────────────────────────────────────────────────────────
%  CHAPTER 10 : CONCLUSION
% ──────────────────────────────────────────────────────────
\chapter{Conclusion}
\label{chap:conclusion}

\section{Bilan du projet}

Ce projet a développé un système complet de prédiction du risque d'incendie de forêt
dans le bassin méditerranéen, depuis l'acquisition et le traitement des données jusqu'à
une interface web interactive. Les contributions principales sont :

\begin{enumerate}[leftmargin=*,itemsep=6pt]
  \item \textbf{Pipeline de données multi-sources} intégrant ERA5 (5 variables
    atmosphériques), MODIS NDVI (état de la végétation) et NASA FIRMS (détections
    actives de feux) sur une grille régulière $0.25°$ couvrant 6\,517 cellules du
    bassin méditerranéen ;

  \item \textbf{Deux architectures de deep learning} — LSTM bi-couche et TCN à
    dilations exponentielles — entraînées avec une perte Weighted MSE adaptée au
    déséquilibre naturel entre cellules à risque et cellules sûres, et évaluées
    sur un test set temporellement disjoint (2023) ;

  \item \textbf{Protocole d'évaluation rigoureux} avec split temporel strict, absence
    de toute fuite d'information entre train et test, et scaler fitté exclusivement
    sur les données d'entraînement ;

  \item \textbf{Interface web interactive} permettant la prédiction en temps réel pour
    tout point géographique de la zone méditerranéenne, la génération de cartes de
    risque sur l'ensemble de la grille, et la comparaison visuelle des performances
    des deux modèles ;

  \item \textbf{Validation sur événements historiques documentés} : les incendies
    de Gironde 2022, d'Attique 2023, d'Evros 2023 et d'Attique 2024 sont correctement
    identifiés comme Élevé ou Critique par les deux modèles.
\end{enumerate}

\section{Apports méthodologiques}

D'un point de vue méthodologique, deux contributions sont particulièrement notables :

\begin{itemize}[leftmargin=*,itemsep=4pt]
  \item La \textbf{construction du risk score} par convolution spatiale sur les
    hotspots futurs + normalisation P99 constitue une formulation élégante du
    problème qui préserve à la fois la continuité spatiale et la robustesse aux
    outliers extrêmes ;
  \item L'approche \textbf{DOY-proxy pour le NDVI} (réutilisation de la climatologie
    2022 pour les années suivantes par correspondance de jour calendaire) est
    scientifiquement fondée pour la région méditerranéenne et compatible avec les
    pratiques des systèmes opérationnels comme EFFIS.
\end{itemize}

\section{Perspectives}

Le système FireRisk constitue une base solide sur laquelle plusieurs extensions
peuvent être construites : intégration de modèles de prévision météorologique NWP
(Numerical Weather Prediction) pour étendre l'horizon de prédiction au-delà de la
réanalyse, couplage avec des modèles de propagation du feu (FARSITE, Prometheus) pour
la gestion opérationnelle des ressources, et extension géographique à d'autres
régions à risque (Californie, Australie, Sibérie).

La lutte contre les incendies de forêt dans le contexte du changement climatique
requiert des outils d'anticipation de plus en plus sophistiqués. Les systèmes fondés
sur le deep learning, capables d'intégrer automatiquement des données hétérogènes
et d'apprendre des patterns complexes, représentent une voie prometteuse vers une
prédiction plus précoce et plus fiable du risque incendie.

\clearpage

% ──────────────────────────────────────────────────────────
%  APPENDICES
% ──────────────────────────────────────────────────────────
\appendix

\chapter{Hyperparamètres complets des modèles}

\begin{center}
\begin{tabular}{lll}
\toprule
\rowcolor{fireOrange!15}
\textbf{Hyperparamètre} & \textcolor{lstmBlue}{\textbf{LSTM}} & \textcolor{tcnRed}{\textbf{TCN}} \\
\midrule
Longueur séquence ($L$)   & 15      & 15 \\
Nombre de features ($F$)  & 8       & 8 \\
Taille cachée ($H$)       & 64      & 64 (canaux) \\
Nombre de couches/blocs   & 2       & 4 \\
Taille noyau              & —       & 3 \\
Dilatations               & —       & $\{1, 2, 4, 8\}$ \\
Dropout                   & 0.2     & 0.2 \\
Tête dense                & 64→1    & 32→1 \\
Activation sortie         & Sigmoid & Sigmoid \\
\midrule
Optimiseur                & Adam    & Adam \\
Learning rate             & $10^{-3}$ & $10^{-3}$ \\
Scheduler                 & ReduceLROnPlateau & ReduceLROnPlateau \\
Facteur scheduler         & 0.5     & 0.5 \\
Patience scheduler        & 5       & 5 \\
Taille batch              & 512     & 512 \\
Epochs max                & 100     & 100 \\
Early stopping patience   & 10      & 10 \\
Poids Weighted MSE        & 10 (si $y > 0.5$) & 10 (si $y > 0.5$) \\
\bottomrule
\end{tabular}
\end{center}

\chapter{Glossaire}

\begin{description}[leftmargin=*,style=nextline]
  \item[\textbf{AUC-ROC}] Area Under the Receiver Operating Characteristic Curve.
    Mesure la capacité discriminante d'un classifieur binaire. $= 1$ : discrimination parfaite ;
    $= 0.5$ : aléatoire.
  \item[\textbf{DOY}] Day of Year — numéro du jour dans l'année (1 = 1er janvier, 365 = 31 décembre).
  \item[\textbf{Early Stopping}] Arrêt de l'entraînement lorsque la performance de
    validation ne s'améliore plus après un nombre donné d'époques (patience).
  \item[\textbf{ERA5}] Cinquième génération de réanalyse atmosphérique de l'ECMWF.
    Couvre 1940–présent, résolution $0.25°$, données horaires.
  \item[\textbf{FIRMS}] Fire Information for Resource Management System (NASA).
    Fournit des détections thermiques actives de feux en quasi-temps-réel.
  \item[\textbf{FRP}] Fire Radiative Power. Puissance radiative du feu mesurée par
    satellite, en mégawatts. Indicateur de l'intensité d'un incendie actif.
  \item[\textbf{LSTM}] Long Short-Term Memory. Architecture de réseau de neurones
    récurrent avec mécanismes de portes pour apprendre des dépendances à long terme.
  \item[\textbf{MAE}] Mean Absolute Error. Erreur absolue moyenne entre prédictions
    et valeurs réelles.
  \item[\textbf{MODIS}] Moderate Resolution Imaging Spectroradiometer. Instrument
    embarqué sur les satellites Terra et Aqua de la NASA.
  \item[\textbf{NDVI}] Normalized Difference Vegetation Index. Indice mesurant la
    verdeur de la végétation depuis l'espace, dans $[-1, 1]$.
  \item[\textbf{NWP}] Numerical Weather Prediction. Prévision numérique du temps par
    simulation des équations de la dynamique atmosphérique.
  \item[\textbf{RMSE}] Root Mean Squared Error. Racine de l'erreur quadratique moyenne.
    Plus sensible aux grandes erreurs que le MAE.
  \item[\textbf{Spearman $\rho$}] Corrélation de rang de Spearman. Mesure la monotonie
    de la relation entre prédiction et valeur réelle, dans $[-1, 1]$.
  \item[\textbf{TCN}] Temporal Convolutional Network. Architecture convolutive causale
    avec dilatations exponentielles pour la modélisation de séquences temporelles.
  \item[\textbf{Weighted MSE}] Mean Squared Error pondérée. Variante de la perte MSE
    qui attribue un poids supérieur aux exemples appartenant à des classes minoritaires.
\end{description}

\chapter{Dépendances logicielles}

\begin{center}
\begin{tabular}{lll}
\toprule
\rowcolor{fireOrange!15}
\textbf{Package} & \textbf{Version} & \textbf{Rôle} \\
\midrule
Python        & $\geq 3.10$  & Langage \\
PyTorch       & $\geq 2.0$   & Deep learning, entraînement GPU/CPU \\
Streamlit     & $\geq 1.30$  & Interface web interactive \\
Plotly        & $\geq 5.18$  & Visualisations interactives \\
pandas        & $\geq 2.0$   & Manipulation de données tabulaires \\
numpy         & $\geq 1.24$  & Calcul numérique, arrays \\
xarray        & $\geq 2023$  & Lecture NetCDF ERA5 \\
rasterio      & $\geq 1.3$   & Lecture GeoTIFF MODIS \\
scipy         & $\geq 1.10$  & Convolution spatiale, Spearman \\
scikit-learn  & $\geq 1.3$   & AUC-ROC \\
pyarrow       & $\geq 13.0$  & Format Parquet \\
\bottomrule
\end{tabular}
\end{center}

\chapter{Références des données}

\begin{enumerate}[leftmargin=*]
  \item \textbf{ERA5} : Hersbach, H. et al. (2020). The ERA5 global reanalysis.
    \textit{Quarterly Journal of the Royal Meteorological Society}, 146(730),
    1999–2049. \url{https://doi.org/10.1002/qj.3803}

  \item \textbf{MODIS NDVI} : Didan, K. (2021). \textit{MODIS/Terra Vegetation
    Indices Monthly L3 Global 1km SIN Grid V061} [MOD13A3]. NASA EOSDIS Land
    Processes DAAC. \url{https://doi.org/10.5067/MODIS/MOD13A3.061}

  \item \textbf{NASA FIRMS} : Davies, D.K., et al. (2009). Detection of active
    fires and biomass burning from Landsat-7 ETM+ data. \textit{Remote Sensing
    of Environment}, 113(1), 69–80.
    \url{https://firms.modaps.eosdis.nasa.gov/}

  \item \textbf{LSTM} : Hochreiter, S. \& Schmidhuber, J. (1997). Long Short-Term
    Memory. \textit{Neural Computation}, 9(8), 1735–1780.

  \item \textbf{TCN} : Bai, S., Kolter, J.Z. \& Koltun, V. (2018). An Empirical
    Evaluation of Generic Convolutional and Recurrent Networks for Sequence
    Modeling. \textit{arXiv:1803.01271}.

  \item \textbf{EFFIS} : San-Miguel-Ayanz, J. et al. (2012). Comprehensive
    monitoring of wildfires in Europe: The European Forest Fire Information System
    (EFFIS). \textit{Approaches to Managing Disaster — Assessing Hazards,
    Emergencies and Disaster Impacts}, InTech.

  \item \textbf{Weighted Loss pour déséquilibre} : King, G. \& Zeng, L. (2001).
    Logistic regression in rare events data.
    \textit{Political Analysis}, 9(2), 137–163.
\end{enumerate}

\end{document}
